\section{Exact torsion costs over wild radical coefficient fields}
\label{sec:false-tate}

The repository's
\texttt{exists\_rootOfUnity\_eval\_le\_gaussVal\_add}
chooses a root of unity at which an integral polynomial does not gain
much valuation. It is a quantitative ingredient in the construction of
arithmetic twists, rather than a consequence of RH.
The earlier graph-tower report treats the coefficient field
$\mathbb Q_3(3^{1/3})$ in its supplement \texttt{wild-cubic.tex}
\cite{priorGraph}. We extend that calculation to every odd prime and
every pure degree-$p$ radical with Eisenstein defining polynomial.
The result determines an exact local obstruction, including all strict
endpoints and simultaneous polynomials with different witnesses.
It does not construct the global Galois characters needed for a new
nonabelian RH proof.

\subsection{Statement and examples}
Let $p$ be an odd prime, let $u\in\mathbb Z_p^\times$, and put
\[
 K=\mathbb Q_p(\pi),\qquad \pi^p=pu.
\]
Normalize the valuation on an algebraic closure by $v(p)=1$ and set
$v(0)=+\infty$. For $h\ge1$ write
\[
 N_h=(p-1)p^{h-1},\quad r_h=N_h^{-1},\quad
 Z_h=\{\zeta-1:\operatorname{ord}(\zeta)=p^h\}.
\]
Fix a finite nonempty set $S$ of positive integers, a radius
$\rho\ge\max_{h\in S}r_h$, and a tolerance $\epsilon\ge0$.
A packet consists of nonzero integral polynomials $F_i\in\mathcal O_K[T]$
and witnesses $x_i$ with $v(x_i)\ge\rho$ and
$t_i=v(F_i(x_i))<\infty$. Its cost and loss at $z$ are
\[
 B=\sum_i t_i,\qquad
 \mathcal L(z)=\sum_i\max\{0,v(F_i(z))-t_i\}.
\]
Different polynomials may have different witnesses. A packet blocks $S$
if $\mathcal L(z)>\epsilon$ for every $z\in\bigcup_{h\in S}Z_h$.
Let $C_K(S,\epsilon)$ be the least cost of a blocking packet.

\begin{theorem}[Exact wild-radical obstruction]\label{thm:false-tate-cost}
For $h\ge2$ put
\[
 \delta_{p,h}=\frac1{p-1}-\frac{2p-1}{pN_h}>0,
 \qquad
 c_{p,h}(\epsilon)=
 \min\left\{1,\frac{\lfloor\epsilon/\delta_{p,h}\rfloor+1}{p}\right\},
\]
and put $c_{p,1}(\epsilon)=1$. Then
\begin{equation}\label{eq:false-tate-cost}
 \boxed{\ C_K(S,\epsilon)=\sum_{h\in S}c_{p,h}(\epsilon).\ }
\end{equation}
The minimum is attained by polynomials with witness zero. Consequently,
if a packet has cost strictly below the right side, some allowed torsion
point satisfies $\mathcal L(z)\le\epsilon$. The strict budget inequality
is necessary. The answer is independent of the unit $u$ and of $\rho$
within the stated range. It remains unchanged for bounded integral
formal power series $F_i\in\mathcal O_K[[T]]$.
\end{theorem}

For one shell the cost is $k/p$ when
$(k-1)\delta_{p,h}\le\epsilon<k\delta_{p,h}$, for
$1\le k\le p-1$, and is $1$ when
$\epsilon\ge(p-1)\delta_{p,h}$. In particular,
$\delta_{p,2}=(p-1)/p^2$. At $p=5$, $h=2$, this is the explicit staircase
\[
\begin{array}{c|ccccc}
\epsilon &[0,4/25)&[4/25,8/25)&[8/25,12/25)&[12/25,16/25)&[16/25,\infty)\\\hline
C_K(\{2\},\epsilon)&1/5&2/5&3/5&4/5&1
\end{array}
\]
At $p=3$ the theorem recovers the predecessor's thresholds
$\delta_{3,h}=1/2-5/(3N_h)$. Thus that special case is inherited;
the general odd-prime and unit-uniform statement is the extension proved
here. For any fixed $p$ the whole tolerance profile changes at every
conductor because $\delta_{p,h}$ increases strictly to $1/(p-1)$.
For each fixed $\epsilon$, however, its value is eventually
\[
 \min\left\{1,\frac{\lfloor(p-1)\epsilon\rfloor+1}{p}\right\}.
\]

\subsection{The degree-exact approximation theorem}
The main arithmetic step determines every monic degree, not merely the
first degree at which a good approximation exists.

\begin{theorem}[Best monic approximation]\label{thm:false-tate-approx}
Fix $h\ge2$, put $N=N_h$, $n=N/p$, and choose $\alpha\in Z_h$.
For every $0\le m<N$,
\begin{equation}\label{eq:false-tate-approx}
 \max_{\substack{P\in\mathcal O_K[T]\text{ monic}\\\deg P=m}}
 v(P(\alpha))
 =\frac mN+\left\lfloor\frac mn\right\rfloor\delta_{p,h}.
\end{equation}
One monic polynomial $B$ of degree $n$ attains all these maxima:
for $m=kn+j$, $0\le j<n$, $0\le k<p$, use $T^jB(T)^k$.
Every root of $B$ has valuation $1/N$, and
\[
 v(B(0))=\frac1p,\qquad
 v(B(\alpha))=\frac1p+\delta_{p,h}
 =\frac{(2p-1)(p^{h-1}-1)}{pN}.
\]
The valuations are the same at every $\alpha\in Z_h$.
For $h=1$, the maximum in degree $0\le m<N_1$ is simply $m/N_1$.
\end{theorem}

\begin{proof}
We give the ramification and integral-basis calculation explicitly.
These are classical local-field tools; related radical-cyclotomic
ramification calculations appear in Viviani \cite{Viviani}, and explicit
uniformizers for $u=1$ in Wang--Yuan \cite{WangYuan}.

Write $C_j=\mathbb Q_p(\zeta_{p^j})$ and $E_j=KC_j$.
The extension $K/\mathbb Q_p$ is totally ramified of degree $p$, with
integer ring $\mathbb Z_p[\pi]$ and different exponent $2p-1$,
since the derivative at $\pi$ is $p\pi^{p-1}$.
It is not Galois: normality would put $\zeta_p$ in a field of degree $p$,
whereas $[\mathbb Q_p(\zeta_p):\mathbb Q_p]=p-1$.

The degree-$p$ extension $E_1/C_1$ is cyclic. The coprime ramification
indices $p$ and $p-1$ show that $E_1/\mathbb Q_p$ is totally ramified
of degree $p(p-1)$. The element
\[
 w=(\zeta_p-1)/\pi
\]
has valuation $1/[p(p-1)]$ and is a uniformizer. A nonidentity element
of $\operatorname{Gal}(E_1/C_1)$ multiplies $w$ by a nontrivial
$p$th root of unity, so the normalized valuation of its difference
from $w$ is $p+1$. Thus the unique lower and upper ramification break
of $E_1/C_1$ is $p$.

The cyclotomic extension $C_h/C_1$ has upper breaks
\[
 p-1,\ 2(p-1),\ldots,(h-1)(p-1).
\]
For example, its lower breaks are $p^k-1$: if
$v_p(a-1)=k$, the automorphism $\zeta\mapsto\zeta^a$ moves the
uniformizer $\zeta-1$ by an element of normalized valuation $p^k$.
Integrating the reciprocal successive indices gives the asserted upper
breaks. The extensions $E_1/C_1$ and $C_h/C_1$ are Galois and linearly
disjoint. Their intersection, of degree dividing $p$, could otherwise
only be $E_1$, which would imply $K\subset C_h$, contrary to
non-Galoisness of $K$ and the abelian nature of $C_h/\mathbb Q_p$.
This proves $[E_h:\mathbb Q_p]=pN$ and $[E_h:K]=N$.

For clarity, total ramification of the compositum is not inferred merely
from total ramification of its two factors. In its abelian group over
$C_1$, a character from the $E_1$ factor has break $p$; nontrivial
characters from the other factor have breaks that are multiples of
$p-1$. These never coincide for odd $p$. The product of two characters
with unequal breaks has their larger break. Consequently no nontrivial
character of the compositum is unramified, proving total ramification
and the product rule for its upper filtration.

For $h\ge2$ the upper-to-lower map sends $p$ to
\[
 (p-1)+p\bigl(p-(p-1)\bigr)=2p-1.
\]
Here the group index is one up to $p-1$ and is $p$ from $p-1$ to $p$.
Lower ramification groups intersect with subgroups, so the cyclic
extension $E_h/C_h$ has its sole lower break at $2p-1$.
Its different exponent is therefore $2p(p-1)$.

Let $d_C$ denote the different exponent of $C_h/\mathbb Q_p$.
The shifted cyclotomic polynomial is Eisenstein over $\mathbb Q_p$,
and $\mathbb Z_p[\alpha]=\mathcal O_{C_h}$. Its degree remains $N$
over $K$, so the discriminant of the power-basis order
$\mathcal O_K[\alpha]$ has valuation $p d_C$ in $K$.
Transitivity of the different through $C_h$ and through $K$ gives
\[
 d(E_h/K)=2p(p-1)+p d_C-(2p-1)N.
\]
All the relevant residue degrees are one. The discriminant-index
formula consequently yields
\begin{equation}\label{eq:false-tate-index}
 I:=\operatorname{length}_{\mathcal O_K}
     (\mathcal O_{E_h}/\mathcal O_K[\alpha])
 =\frac{(2p-1)N-2p(p-1)}2.
\end{equation}

Choose a monic degree-$n$ polynomial $B$ maximizing
$b=v(B(\alpha))$. Such a maximum exists: its lower coefficients range
in the compact set $\mathcal O_K^n$, and the evaluation map has compact
image not containing zero because $\alpha$ has degree $N>n$ over $K$.
The trial polynomial $T^n$ gives $b\ge1/p$.
We claim that $b\notin N^{-1}\mathbb Z$.
If it were, write
\[
 b=j/N+t/p,\qquad 0\le j<n,\quad t\ge1.
\]
Both $B(\alpha)$ and $\pi^t\alpha^j$ have this valuation. Since
$E_h$ and $K$ have residue field $\mathbb F_p$, some unit of
$\mathcal O_K$ cancels their initial residue. Subtracting its multiple
of $\pi^tT^j$ from $B$ gives a better monic integral polynomial of
the same degree, a contradiction.

The $N$ elements
\[
 w_{jk}=\alpha^jB(\alpha)^k,
 \qquad 0\le j<n,\quad 0\le k<p,
\]
form an $\mathcal O_K$-basis of $\mathcal O_K[\alpha]$, by triangularity
in their distinct degrees $kn+j$. Their valuations are pairwise
incongruent modulo $p^{-1}\mathbb Z$. Indeed $b$ has order $p$ in
$(pN)^{-1}\mathbb Z/N^{-1}\mathbb Z$, so equality of two such classes
first forces equality of $k$ and then of $j$.
Every nonzero $K$-linear combination therefore has valuation equal
to its uniquely least term.

It follows that
$w_{jk}/\pi^{\lfloor p v(w_{jk})\rfloor}$ form an integral basis
of $E_h/K$: their valuations are distinct in $[0,1/p)$, so a linear
combination is integral exactly when all its coefficients are integral.
Put $\delta=b-1/p$. The valuations in degree $m=kn+j$ are
$m/N+k\delta$. Their fractional parts after multiplication by $p$
are precisely $0,1/N,\ldots,(N-1)/N$. Thus
\begin{align*}
 I&=\sum_{j,k}\lfloor p v(w_{jk})\rfloor\\
  &=\frac{(p-1)(N-1)}2+
       \frac{pN(p-1)}2\,\delta.
\end{align*}
Comparison with \eqref{eq:false-tate-index} gives
\[
 \delta=\frac1{p-1}-\frac{2p-1}{pN}.
\]
In the monic triangular basis $T^jB^k$, every monic degree-$m$
polynomial has coefficient one at its leading basis element.
The no-cancellation property bounds its valuation by that leading
basis value, and the basis element itself attains it. This proves
\eqref{eq:false-tate-approx}.

We still need the assertion about the roots of $B$, because it is
what makes the obstruction cost exactly $1/p$. Collect its roots of
valuation $r=1/N$ into a monic factor $P$ of degree $a$. The factor
lies in $K[T]$ since the valuation is Galois invariant, and is integral.
Every other root contributes at most $r$ to $v(B(\alpha))$.
If $a<n$, formula \eqref{eq:false-tate-approx} in degree $a$ gives
$v(P(\alpha))\le a/N$. It would follow that
$v(B(\alpha))\le n/N$, contrary to $\delta>0$.
Therefore all roots have valuation $r$, and $v(B(0))=n/N=1/p$.
The equality of valuations over the whole shell follows from
$\operatorname{Gal}(E_h/K)$ acting transitively on the conjugates of
$\alpha$ and preserving the valuation.

Finally, at $h=1$ the values $v(a_j)+j/(p-1)$ of distinct nonzero
coefficient terms of degree below $p-1$ are incongruent modulo
$p^{-1}\mathbb Z$, since $p$ and $p-1$ are coprime. A monic polynomial
of degree $m<p-1$ therefore cannot exceed $m/(p-1)$ in valuation.
The monomial attains it.
\end{proof}

\subsection{From monic approximation to arbitrary packets}
\begin{proof}[Proof of Theorem~\ref{thm:false-tate-cost}]
For $F=\sum a_jT^j$ let
$G_s(F)=\min_j\{v(a_j)+js\}$. Factoring over an algebraic closure gives
\[
 G_s(F)=v(a_{\deg F})+\sum_{F(\beta)=0}\min\{s,v(\beta)\},
\]
with multiplicities. Fix a shell of radius $r=1/N$, and let $P_i$
collect exactly the roots of $F_i$ having valuation $r$.
These monic factors are integral and defined over $K$.
Because $t_i\ge G_\rho(F_i)\ge G_r(F_i)$, and because nonmatching roots
have $v(z-\beta)=\min\{r,v(\beta)\}$, we have
\[
 \max\{0,v(F_i(z))-t_i\}
 \le\sum_{P_i(\beta)=0}\bigl(v(z-\beta)-r\bigr).
\]
Every summand on the right is nonnegative. For $P=\prod_iP_i$,
$m=\deg P$, this proves
\begin{equation}\label{eq:false-tate-loss}
 \mathcal L(z)\le v(P(z))-m/N.
\end{equation}
Moreover, matching roots contribute $m/N$ to the sum of
$G_\rho(F_i)-G_0(F_i)$, and $G_0(F_i)\ge0$; hence the packet's cost
is at least $m/N$.

For $m<N$, Theorem~\ref{thm:false-tate-approx} bounds the right side
of \eqref{eq:false-tate-loss} by $\lfloor m/n\rfloor\delta_{p,h}$.
The least degree that can make this exceed $\epsilon$ is $kn$, where
$k=\lfloor\epsilon/\delta_{p,h}\rfloor+1$, provided $k<p$.
It costs $k/p$. At degree $N$ the shifted cyclotomic polynomial
vanishes on the whole shell and has constant term of valuation one.
For $k<p$, the polynomial $B^k$ with witness zero has cost $k/p$
and loss exactly $k\delta_{p,h}$. This proves the single-shell formula
with its strict endpoints. The $h=1$ case follows from the last part
of the approximation theorem and the same cyclotomic construction.

For several shells, their matching-root multisets are disjoint.
Applying the preceding lower bound separately and summing their
contributions to $\sum_i(G_\rho(F_i)-G_0(F_i))$ proves the lower
bound by the sum in \eqref{eq:false-tate-cost}.
For attainment, take one of the optimal blocking polynomials for each
shell as separate members of the packet, all with witness zero.
Each polynomial has every root on its own shell. On another shell its
valuation minus its constant-term valuation is nonpositive, so it
contributes zero to the positive-part loss. On its own shell it gives
the required strict loss. The costs therefore add exactly.

For bounded integral power series, truncate to a sufficiently high degree
that the omitted terms have valuation greater than
$\sum_i t_i+\epsilon$ at every witness and every allowed torsion point.
This is possible because all these arguments have valuation bounded
below by a fixed positive number. The witness values remain unchanged.
Each loss at most $\epsilon$ is unchanged, and each loss exceeding
$\epsilon$ still exceeds it, so blocking is preserved. The polynomial
lower bound applies, and polynomial examples already attain it.
\end{proof}

\subsection{Independent exact arithmetic and scope}
The accompanying \texttt{check\_false\_tate.py} constructs a maximizing
polynomial by successive residue cancellation for $u=1$. It evaluates
$P(\zeta-1)$ using
\[
 \operatorname{Res}_{Y}\bigl(Y^p-p,\ P(T)\bigr)
       \pmod{\Phi_{p^h}(1+T)},
\]
where the coefficients of $P$ are written as polynomials in $Y$.
If this remainder is $\sum_{j=0}^{N-1}c_jT^j$, then
\[
 pN\,v(P(\alpha))=\min_{c_j\ne0}\{Nv_p(c_j)+j\}.
\]
The formula uses the Eisenstein power basis in $C_h$ and the degree-$p$
norm from $E_h$; no discriminant or ramification-break formula enters
the calculation. The checked values are
\[
\begin{array}{c|r|r|c|r}
(p,h)&N&\deg B&\delta_{p,h}&pN\,v(B(\alpha))\\\hline
(3,2)&6&2&2/9&10\\
(3,3)&18&6&11/27&40\\
(3,4)&54&18&38/81&130\\
(5,2)&20&4&4/25&36\\
(5,3)&100&20&29/125&216\\
(7,2)&42&6&6/49&78
\end{array}
\]
The JSON certificate records every coefficient and cancellation step,
and verifies that the resulting monic basis has distinct valuation
classes. These are exact finite checks, not a proof by experimental
extrapolation; the all-prime and all-conductor proof is above.

The radical-cyclotomic extensions themselves and their ramification
are classical \cite{Viviani,WangYuan}. The integral-basis viewpoint
belongs to the established theory of higher Newton polygons and
reduced bases \cite{GMN}. The contribution relative to the supplied
reports is the general formula for the strict simultaneous-loss
optimization over all odd-prime radical coefficient fields.
A search of the cited primary literature did not establish publication
priority for that optimization, so it should not be described as a
new named conjecture being settled. Its practical use is an exact
local budget for candidate twist families; global realization,
Selmer comparison, and prime-sum estimates remain separate problems.
